% NSF NCAR / UCAR beamer theme -- math font specimen.
%
% One slide per mathfont= preset, and a summary slide with every preset on
% it (docs/math-fonts.png).  \usetheme[mathfont=<preset>]{NCAR} picks one for
% a deck; here \ncarmathversion sets each up as a math version so a single
% document can show them all (poppins, which cannot be a version, is the
% document's normal math).  Build: make examples/mathfonts.pdf
%
\documentclass[aspectratio=169]{beamer}
\usetheme[mathfont=keep,logo=false]{NCAR}

\ncarmathversion{serif}{serif}
\ncarmathversion{lm}{lm}
\ncarmathversion{stix}{stix}
\ncarmathversion{pagella}{pagella}
\ncarmathversion{sans}{sans}
\ncarmathversion{fira}{fira}
\ncarmathfont{poppins}

\title{Math fonts}
\subtitle{\texttt{mathfont=} presets of the NCAR theme}
\author{\texttt{\textbackslash usetheme[mathfont=serif]\{NCAR\}} (the default)}
\date{}

% the same equations under each version
\newcommand\specimen{%
  \begin{equation*}
    \frac{\partial \mathbf{u}}{\partial t} + (\mathbf{u}\cdot\nabla)\mathbf{u}
      = -\frac{1}{\rho}\nabla p + \nu\nabla^2\mathbf{u} + \mathbf{g},
    \qquad
    \int_{-\infty}^{\infty} e^{-x^2}\,\mathrm{d}x = \sqrt{\pi},
    \qquad
    \sum_{n=1}^{\infty}\frac{1}{n^2} = \frac{\pi^2}{6}
  \end{equation*}
  \begin{equation*}
    \hat{f}(\xi) = \int_{\mathbb{R}^n} f(x)\, e^{-2\pi i\, x\cdot\xi}\,\mathrm{d}x,
    \qquad
    \mathcal{L}\{y\}(s) = \int_0^\infty y(t)\,e^{-st}\,\mathrm{d}t,
    \qquad
    \mathbf{A} = \begin{pmatrix} \alpha_{11} & \alpha_{12} \\ \alpha_{21} & \Omega_{22} \end{pmatrix}
  \end{equation*}
  where $\mathbf{u}\in\mathbb{R}^3$ is the velocity, $\rho$ the density and
  $\nu$ the kinematic viscosity; $\Gamma(z) = \int_0^\infty t^{z-1}e^{-t}\,\mathrm{d}t$
  for $\operatorname{Re}(z) > 0$, and $\epsilon \ll \delta \ll 1$.}

% \specimenframe[<math version>]{<preset>}{<title>}{<what it is>}
\newcommand\specimenframe[4][]{%
  \begin{frame}{\texttt{mathfont=#2}}{#3}
    \mathversion{\ifx\relax#1\relax#2\else#1\fi}
    \specimen
    \vfill
    {\small\color{DarkBlue!70}#4}
  \end{frame}}

% \specimenrow[<math version>]{<preset>}: one line of the summary slide
\newcommand\specimenrow[2][]{%
  \par\noindent\makebox[\textwidth][l]{\small
    \makebox[0.1\textwidth][l]{\texttt{#2}}%
    \mathversion{\ifx\relax#1\relax#2\else#1\fi}%
    $\dfrac{\partial \mathbf{u}}{\partial t} + (\mathbf{u}\cdot\nabla)\mathbf{u}
      = -\dfrac{1}{\rho}\nabla p + \nu\nabla^2\mathbf{u}$\quad
    $\int_0^{\infty} e^{-x^2}\,\mathrm{d}x = \tfrac{1}{2}\sqrt{\pi}$\quad
    $\hat{f}(\xi),\ \mathbb{R}^n,\ \Gamma\Omega\,\alpha\beta\gamma$\quad
    text $\rho$}%
  \par\vspace{0.1ex}}

\begin{document}

\begin{frame}[plain,noframenumbering]
  \titlepage
\end{frame}

\specimenframe{serif}{New Computer Modern Math, Book weight}{The default.
  The Computer Modern design at a weight that holds its own beside Poppins;
  Latin Modern Math stands in where it is not installed.}
\specimenframe{lm}{Latin Modern Math}{The lighter Computer Modern most LaTeX
  users know; what Quarto PDFs got before theme 2.3.}
\specimenframe{stix}{STIX Two Math}{Times-like; dense equations stay compact.}
\specimenframe{pagella}{TeX Gyre Pagella Math}{Palatino-like; the widest
  letterforms of the serif presets.}
\specimenframe{sans}{Lete Sans Math (Lato)}{A complete sans math font, the
  nearest relative to Poppins; Fira Math stands in where it is not installed.}
\specimenframe{fira}{Fira Math}{Humanist sans with a large x-height; it
  has no bold or script alphabets (the box above is \texttt{\textbackslash mathcal\{L\}}).}
\specimenframe[normal]{poppins}{Poppins letters over Computer Modern symbols}{The look
  of theme versions before 2.3: beamer's own sans math.  Fine for a variable
  name in a sentence, less so for anything with fractions or operators.}

\begin{frame}{The presets side by side}
  \specimenrow{serif}\specimenrow{lm}\specimenrow{stix}\specimenrow{pagella}
  \specimenrow{sans}\specimenrow{fira}\specimenrow[normal]{poppins}
\end{frame}

\end{document}
